% ==============================================================================
% CHAPTER 9: PRACTICAL DEEP LEARNING
% ==============================================================================

\chapter{Practical Deep Learning}
\label{ch:practical_dl}

\begin{keyconcept}
Moving from theory to implementation requires understanding \textbf{automatic differentiation} (autograd), which is the computational engine behind modern deep learning frameworks. This chapter bridges the gap between mathematical understanding and practical implementation.
\end{keyconcept}

% ------------------------------------------------------------------------------
\section{From Theory to Implementation}
% ------------------------------------------------------------------------------

\subsection{Three Types of Differentiation}

We've derived backpropagation mathematically. Now we need to implement it computationally. There are three approaches:

\begin{enumerate}
    \item \textbf{Symbolic Differentiation}: Manipulate mathematical expressions
    \begin{itemize}
        \item Use symbols and calculus rules to derive exact derivatives
        \item Example: $\frac{d}{dx}(x^2) = 2x$
        \item Pro: Exact results
        \item Con: Expression explosion for complex functions; inefficient
    \end{itemize}
    
    \item \textbf{Numerical Differentiation}: Approximate using finite differences
    \begin{itemize}
        \item Compute $f'(x) \approx \frac{f(x + \epsilon) - f(x)}{\epsilon}$
        \item Pro: Simple to implement
        \item Con: Slow (requires many function evaluations); numerically unstable
    \end{itemize}
    
    \item \textbf{Automatic Differentiation (Autograd)}: Best of both worlds
    \begin{itemize}
        \item Computes exact derivatives efficiently
        \item Breaks computation into elementary operations
        \item Applies chain rule automatically
        \item This is what PyTorch, TensorFlow, JAX use!
    \end{itemize}
\end{enumerate}

\subsection{Numerical Differentiation Details}

\textbf{Forward Difference:}
\[
f'(x) \approx \frac{f(x + \epsilon) - f(x)}{\epsilon}
\]

\textbf{Central Difference (more accurate):}
\[
f'(x) \approx \frac{f(x + \epsilon) - f(x - \epsilon)}{2\epsilon}
\]

\begin{lstlisting}[language=Python, caption={Numerical Differentiation}]
def numerical_gradient(f, x, epsilon=1e-5):
    """Compute gradient numerically using central difference."""
    grad = np.zeros_like(x)
    for i in range(len(x)):
        x_plus = x.copy()
        x_minus = x.copy()
        x_plus[i] += epsilon
        x_minus[i] -= epsilon
        grad[i] = (f(x_plus) - f(x_minus)) / (2 * epsilon)
    return grad
\end{lstlisting}

\textbf{Problem}: For $n$ parameters, this requires $2n$ function evaluations. With millions of parameters, this is far too slow!

% ------------------------------------------------------------------------------
\section{Automatic Differentiation (Autograd)}
% ------------------------------------------------------------------------------

\subsection{The Core Idea}

Autograd works by:
\begin{enumerate}
    \item \textbf{Recording} operations in a computational graph during the forward pass
    \item \textbf{Traversing} the graph backward to compute gradients
    \item Each operation stores its \textbf{local gradient} and how to compute it
\end{enumerate}

\subsection{Forward Mode vs. Reverse Mode}

\begin{itemize}
    \item \textbf{Forward Mode}: Computes derivatives of all outputs w.r.t. \textit{one} input
    \begin{itemize}
        \item Good when: many outputs, few inputs
        \item Complexity: $O(\text{inputs})$ passes
    \end{itemize}
    
    \item \textbf{Reverse Mode}: Computes derivatives of \textit{one} output w.r.t. all inputs
    \begin{itemize}
        \item Good when: one output (loss), many inputs (parameters)
        \item Complexity: $O(1)$ pass for all gradients
        \item \textbf{This is what backpropagation uses!}
    \end{itemize}
\end{itemize}

\begin{important}
Reverse-mode automatic differentiation is \textbf{exactly} backpropagation. When you call \texttt{loss.backward()} in PyTorch, you're running reverse-mode autograd.
\end{important}

% ------------------------------------------------------------------------------
\section{Building a Simple Autograd System}
% ------------------------------------------------------------------------------

Let's build a minimal autograd engine to understand how frameworks like PyTorch work.

\subsection{The Value Class}

\begin{lstlisting}[language=Python, caption={Simple Autograd Engine: Value Class}]
class Value:
    """A scalar value that tracks gradients."""
    
    def __init__(self, data, _children=(), _op=''):
        self.data = data
        self.grad = 0.0  # Gradient w.r.t. this value
        
        # For building the computational graph
        self._backward = lambda: None  # How to compute parent gradients
        self._prev = set(_children)    # Parent nodes
        self._op = _op                 # Operation that created this
    
    def __add__(self, other):
        other = other if isinstance(other, Value) else Value(other)
        out = Value(self.data + other.data, (self, other), '+')
        
        def _backward():
            # Gradient of a+b w.r.t. a is 1, w.r.t. b is 1
            self.grad += 1.0 * out.grad
            other.grad += 1.0 * out.grad
        out._backward = _backward
        
        return out
    
    def __mul__(self, other):
        other = other if isinstance(other, Value) else Value(other)
        out = Value(self.data * other.data, (self, other), '*')
        
        def _backward():
            # Gradient of a*b w.r.t. a is b, w.r.t. b is a
            self.grad += other.data * out.grad
            other.grad += self.data * out.grad
        out._backward = _backward
        
        return out
    
    def __pow__(self, n):
        out = Value(self.data ** n, (self,), f'**{n}')
        
        def _backward():
            # Gradient of x^n w.r.t. x is n*x^(n-1)
            self.grad += n * (self.data ** (n - 1)) * out.grad
        out._backward = _backward
        
        return out
    
    def tanh(self):
        x = self.data
        t = (math.exp(2*x) - 1) / (math.exp(2*x) + 1)
        out = Value(t, (self,), 'tanh')
        
        def _backward():
            # Gradient of tanh(x) is 1 - tanh(x)^2
            self.grad += (1 - t**2) * out.grad
        out._backward = _backward
        
        return out
    
    def backward(self):
        """Run backpropagation from this node."""
        # Topological sort
        topo = []
        visited = set()
        def build_topo(v):
            if v not in visited:
                visited.add(v)
                for child in v._prev:
                    build_topo(child)
                topo.append(v)
        build_topo(self)
        
        # Backprop
        self.grad = 1.0  # dL/dL = 1
        for node in reversed(topo):
            node._backward()
    
    def __repr__(self):
        return f"Value(data={self.data:.4f}, grad={self.grad:.4f})"
    
    # Additional operations
    def __radd__(self, other): return self + other
    def __rmul__(self, other): return self * other
    def __neg__(self): return self * -1
    def __sub__(self, other): return self + (-other)
\end{lstlisting}

\subsection{Using Our Autograd Engine}

\begin{lstlisting}[language=Python, caption={Example: Computing Gradients}]
# Define values
a = Value(2.0)
b = Value(-3.0)
c = Value(10.0)

# Build expression: d = a*b + c
d = a * b + c

# Backward pass
d.backward()

# Check gradients
print(f"a.grad = {a.grad}")  # dd/da = b = -3
print(f"b.grad = {b.grad}")  # dd/db = a = 2
print(f"c.grad = {c.grad}")  # dd/dc = 1
\end{lstlisting}

\textbf{Output:}
\begin{verbatim}
a.grad = -3.0
b.grad = 2.0
c.grad = 1.0
\end{verbatim}

% ------------------------------------------------------------------------------
\subsection{Visualizing the Computation Graph}
% ------------------------------------------------------------------------------

As expressions become complex, visualizing the underlying DAG is incredibly useful. We can write a topological tracing function and use the \texttt{graphviz} library to render the graph.

\begin{lstlisting}[language=Python, caption={Tracing and Rendering with Graphviz}]
from graphviz import Digraph

def trace(root):
    # Builds a set of all nodes and edges in a graph
    nodes, edges = set(), set()
    def build(v):
        if v not in nodes:
            nodes.add(v)
            for child in v._prev:
                edges.add((child, v))
                build(child)
    build(root)
    return nodes, edges

def draw_dot(root):
    dot = Digraph(format='svg', graph_attr={'rankdir': 'LR'}) # Left to Right
    
    nodes, edges = trace(root)
    for n in nodes:
        uid = str(id(n))
        # Create a rectangular record node for the value
        dot.node(name=uid, label=f"{{ data {n.data:.4f} | grad {n.grad:.4f} }}", shape='record')
        if n._op:
            # If this value is a result of some operation, create an op node for it
            dot.node(name=uid + n._op, label=n._op)
            # and connect this node to it
            dot.edge(uid + n._op, uid)
    
    for n1, n2 in edges:
        # Connect n1 to the op node of n2
        dot.edge(str(id(n1)), str(id(n2)) + n2._op)
        
    return dot

# Visualize the graph
# draw_dot(d)  # Renders the SVG graph
\end{lstlisting}

% ------------------------------------------------------------------------------
\section{Building a Neural Network from Scratch}
% ------------------------------------------------------------------------------

\subsection{Neuron Class}

\begin{lstlisting}[language=Python, caption={Neuron Implementation}]
import random

class Neuron:
    """A single neuron with weights, bias, and activation."""
    
    def __init__(self, nin, activation='tanh'):
        # Initialize weights randomly
        self.w = [Value(random.uniform(-1, 1)) for _ in range(nin)]
        self.b = Value(0.0)
        self.activation = activation
    
    def __call__(self, x):
        # Weighted sum: w.x + b
        act = sum((wi * xi for wi, xi in zip(self.w, x)), self.b)
        
        # Apply activation
        if self.activation == 'tanh':
            return act.tanh()
        elif self.activation == 'linear':
            return act
        else:
            raise ValueError(f"Unknown activation: {self.activation}")
    
    def parameters(self):
        return self.w + [self.b]
\end{lstlisting}

\subsection{Layer Class}

\begin{lstlisting}[language=Python, caption={Layer Implementation}]
class Layer:
    """A layer of neurons."""
    
    def __init__(self, nin, nout, activation='tanh'):
        self.neurons = [Neuron(nin, activation) for _ in range(nout)]
    
    def __call__(self, x):
        outputs = [n(x) for n in self.neurons]
        return outputs[0] if len(outputs) == 1 else outputs
    
    def parameters(self):
        return [p for n in self.neurons for p in n.parameters()]
\end{lstlisting}

\subsection{MLP Class}

\begin{lstlisting}[language=Python, caption={Multi-Layer Perceptron}]
class MLP:
    """Multi-Layer Perceptron."""
    
    def __init__(self, nin, layer_sizes):
        sizes = [nin] + layer_sizes
        self.layers = []
        
        for i in range(len(layer_sizes)):
            # Use 'linear' for output layer, 'tanh' for hidden
            activation = 'linear' if i == len(layer_sizes) - 1 else 'tanh'
            self.layers.append(Layer(sizes[i], sizes[i+1], activation))
    
    def __call__(self, x):
        for layer in self.layers:
            x = layer(x)
        return x
    
    def parameters(self):
        return [p for layer in self.layers for p in layer.parameters()]
    
    def zero_grad(self):
        for p in self.parameters():
            p.grad = 0.0
\end{lstlisting}

% ------------------------------------------------------------------------------
\section{The Training Loop}
% ------------------------------------------------------------------------------

\begin{lstlisting}[language=Python, caption={Complete Training Loop}]
# Dataset: simple XOR-like problem
X = [
    [2.0, 3.0, -1.0],
    [3.0, -1.0, 0.5],
    [0.5, 1.0, 1.0],
    [1.0, 1.0, -1.0],
]
y = [1.0, -1.0, -1.0, 1.0]  # Target outputs

# Create network: 3 inputs -> 4 hidden -> 4 hidden -> 1 output
model = MLP(3, [4, 4, 1])

# Training hyperparameters
learning_rate = 0.1
epochs = 100

# Training loop
for epoch in range(epochs):
    # Forward pass: compute predictions
    predictions = [model(x) for x in X]
    
    # Compute loss (MSE)
    loss = sum((pred - target)**2 for pred, target in zip(predictions, y))
    
    # Zero gradients (important!)
    model.zero_grad()
    
    # Backward pass
    loss.backward()
    
    # Gradient descent update
    for p in model.parameters():
        p.data -= learning_rate * p.grad
    
    # Print progress
    if epoch % 10 == 0:
        print(f"Epoch {epoch}, Loss: {loss.data:.4f}")

# Final predictions
print("\nFinal predictions:")
for x, target in zip(X, y):
    pred = model(x)
    print(f"  Input: {x} -> Pred: {pred.data:.3f}, Target: {target}")
\end{lstlisting}

\textbf{Sample Output:}
\begin{verbatim}
Epoch 0, Loss: 4.5231
Epoch 10, Loss: 1.2847
Epoch 20, Loss: 0.3412
Epoch 50, Loss: 0.0128
Epoch 90, Loss: 0.0012

Final predictions:
  Input: [2.0, 3.0, -1.0] -> Pred: 0.989, Target: 1.0
  Input: [3.0, -1.0, 0.5] -> Pred: -0.992, Target: -1.0
  Input: [0.5, 1.0, 1.0] -> Pred: -0.985, Target: -1.0
  Input: [1.0, 1.0, -1.0] -> Pred: 0.991, Target: 1.0
\end{verbatim}

% ------------------------------------------------------------------------------
\section{Training on Complex Datasets}
% ------------------------------------------------------------------------------

While XOR is a fundamental nonlinear problem, real-world data is far messier. We can test our custom framework on classic machine learning datasets using \textbf{mini-batch Stochastic Gradient Descent (SGD)}.

\subsection{The Make Moons Dataset}

The ``Make Moons'' dataset from \texttt{scikit-learn} generates two interleaving half circles, perfect for testing non-linear classification.

\begin{lstlisting}[language=Python, caption={Training on Make Moons with Mini-batches}]
from sklearn.datasets import make_moons
import numpy as np

# Generate dataset
X, y = make_moons(n_samples=100, noise=0.1)
y = y*2 - 1 # Convert labels 0, 1 to -1, 1 for our tanh output

# Initialize model (2 inputs -> 16 hidden -> 16 hidden -> 1 output)
model = MLP(2, [16, 16, 1])

# Mini-batch Training Loop
batch_size = 32
epochs = 100
learning_rate = 0.05

for epoch in range(epochs):
    # Shuffle indices for SGD
    indices = np.random.permutation(len(X))
    
    for i in range(0, len(X), batch_size):
        batch_idx = indices[i:i+batch_size]
        X_batch, y_batch = X[batch_idx], y[batch_idx]
        
        # Forward pass
        preds = [model(x.tolist()) for x in X_batch]
        
        # Mean Squared Error Loss
        loss = sum((pred - target)**2 for pred, target in zip(preds, y_batch)) 
        loss = loss * (1.0/len(y_batch))
        
        # L2 Regularization
        reg_loss = 1e-4 * sum((p*p for p in model.parameters()))
        total_loss = loss + reg_loss
        
        # Backward pass
        model.zero_grad()
        total_loss.backward()
        
        # Update
        for p in model.parameters():
            p.data -= learning_rate * p.grad
\end{lstlisting}

By swapping the dataset generation to \texttt{make\_circles} or using the \textbf{Iris dataset}, the exact same engine can model a variety of complex feature spaces, proving that our from-scratch code functions equivalently to industrial frameworks.

% ------------------------------------------------------------------------------
\section{The Complete Training Algorithm}
% ------------------------------------------------------------------------------

\begin{algorithm}[Neural Network Training]
\textbf{Input:} Dataset $\mathcal{D} = \{(x_i, y_i)\}$, architecture, hyperparameters

\begin{enumerate}
    \item \textbf{Initialize} weights and biases randomly
    \item \textbf{Repeat} for each epoch:
    \begin{enumerate}
        \item \textbf{Shuffle} the dataset
        \item \textbf{For each mini-batch}:
        \begin{enumerate}
            \item \textbf{Forward pass}: Compute predictions $\hat{y} = f_\theta(x)$
            \item \textbf{Compute loss}: $\mathcal{L} = \text{LossFunction}(\hat{y}, y)$
            \item \textbf{Zero gradients}: $\nabla_\theta = 0$
            \item \textbf{Backward pass}: Compute $\nabla_\theta \mathcal{L}$ via autograd
            \item \textbf{Update parameters}: $\theta \leftarrow \theta - \eta \nabla_\theta \mathcal{L}$
        \end{enumerate}
    \end{enumerate}
    \item \textbf{Return} trained parameters $\theta$
\end{enumerate}
\end{algorithm}

\begin{important}
\textbf{Don't forget to zero gradients!} Gradients accumulate by default (useful for certain architectures), so you must explicitly reset them before each backward pass.
\end{important}

% ------------------------------------------------------------------------------
\section{Hyperparameters}
% ------------------------------------------------------------------------------

\subsection{Key Hyperparameters}

\begin{center}
\begin{tabular}{lll}
\toprule
\textbf{Hyperparameter} & \textbf{Typical Range} & \textbf{Notes} \\
\midrule
Learning rate $\eta$ & $10^{-4}$ to $10^{-1}$ & Most important; start with $0.001$ \\
Batch size & 32, 64, 128, 256 & Larger = more stable gradients \\
Number of layers & 2--10+ & Deeper = more representation power \\
Neurons per layer & 32--1024 & Task-dependent \\
Epochs & 10--1000+ & Until validation loss plateaus \\
\bottomrule
\end{tabular}
\end{center}

\subsection{Learning Rate}

The learning rate $\eta$ is arguably the most important hyperparameter:

\begin{itemize}
    \item \textbf{Too high}: Overshooting, unstable training, loss explodes
    \item \textbf{Too low}: Very slow convergence, may get stuck
    \item \textbf{Just right}: Steady decrease in loss
\end{itemize}

\textbf{Learning Rate Schedules:}
\begin{itemize}
    \item \textbf{Step decay}: Reduce by factor every $n$ epochs
    \item \textbf{Exponential decay}: $\eta_t = \eta_0 \cdot \gamma^t$
    \item \textbf{Cosine annealing}: Smooth decrease following cosine curve
    \item \textbf{Warmup}: Start small, increase gradually, then decay
\end{itemize}

% ------------------------------------------------------------------------------
\section{Common Challenges}
% ------------------------------------------------------------------------------

\subsection{Vanishing Gradients}

\textbf{Problem}: In deep networks, gradients become exponentially small as they backpropagate, especially with sigmoid/tanh activations.

\textbf{Why it happens}:
\begin{itemize}
    \item Sigmoid derivative: $\sigma'(x) = \sigma(x)(1 - \sigma(x)) \leq 0.25$
    \item After $L$ layers: $\text{gradient} \propto 0.25^L \to 0$
\end{itemize}

\textbf{Solutions}:
\begin{itemize}
    \item Use ReLU activation: $\text{ReLU}'(x) = 1$ for $x > 0$
    \item Use residual connections (skip connections)
    \item Careful weight initialization (Xavier, He)
    \item Batch normalization
\end{itemize}

\subsection{Exploding Gradients}

\textbf{Problem}: Gradients become extremely large, causing numerical overflow.

\textbf{Solutions}:
\begin{itemize}
    \item \textbf{Gradient clipping}: Cap gradients at maximum value
    \[
    \nabla \leftarrow \min\left(1, \frac{\text{threshold}}{||\nabla||}\right) \cdot \nabla
    \]
    \item Proper weight initialization
    \item Lower learning rate
\end{itemize}

\subsection{Overfitting}

\textbf{Problem}: Model memorizes training data but generalizes poorly.

\textbf{Signs}: Training loss $\ll$ validation loss.

\textbf{Solutions}:
\begin{itemize}
    \item \textbf{More data}: Always the best solution
    \item \textbf{Regularization}: L2 (weight decay), L1, Dropout
    \item \textbf{Early stopping}: Stop when validation loss increases
    \item \textbf{Simpler model}: Fewer layers/neurons
    \item \textbf{Data augmentation}: Artificially expand dataset
\end{itemize}

\subsection{Underfitting}

\textbf{Problem}: Model is too simple to capture patterns.

\textbf{Signs}: Both training and validation loss are high.

\textbf{Solutions}:
\begin{itemize}
    \item More complex model (more layers/neurons)
    \item Train longer
    \item Better features
    \item Reduce regularization
\end{itemize}

% ------------------------------------------------------------------------------
\section{Weight Initialization}
% ------------------------------------------------------------------------------

Proper initialization is crucial for training deep networks.

\subsection{Xavier/Glorot Initialization}

For layers with tanh/sigmoid activation:
\[
W \sim \mathcal{U}\left(-\sqrt{\frac{6}{n_{in} + n_{out}}}, \sqrt{\frac{6}{n_{in} + n_{out}}}\right)
\]

\subsection{He Initialization}

For layers with ReLU activation:
\[
W \sim \mathcal{N}\left(0, \sqrt{\frac{2}{n_{in}}}\right)
\]

\begin{infobox}{Why Initialization Matters}
Poor initialization can cause:
\begin{itemize}
    \item All neurons outputting similar values (dead neurons)
    \item Gradients vanishing or exploding from the start
    \item Slow or impossible training
\end{itemize}
Good initialization keeps activations and gradients in a reasonable range throughout the network.
\end{infobox}

% ------------------------------------------------------------------------------
\section{Regularization Techniques}
% ------------------------------------------------------------------------------

\subsection{L2 Regularization (Weight Decay)}

Add penalty for large weights:
\[
\mathcal{L}_{\text{reg}} = \mathcal{L} + \frac{\lambda}{2} \sum_i w_i^2
\]

This encourages smaller weights, which often leads to better generalization.

\subsection{Dropout}

During training, randomly set neurons to zero with probability $p$:
\begin{itemize}
    \item Prevents co-adaptation of neurons
    \item Acts like training an ensemble of networks
    \item At test time, scale outputs by $(1-p)$ or use inverted dropout
\end{itemize}

% ------------------------------------------------------------------------------
\section{Tips for Success}
% ------------------------------------------------------------------------------

\begin{summary}
\textbf{Practical Guidelines:}

\begin{enumerate}
    \item \textbf{Start simple}: Begin with a small network and gradually increase complexity
    
    \item \textbf{Overfit first}: Make sure your model can memorize a small batch before scaling up
    
    \item \textbf{Monitor everything}: Track training/validation loss, gradients, activations
    
    \item \textbf{Use established architectures}: Don't reinvent the wheel; start with proven designs
    
    \item \textbf{Normalize inputs}: Zero mean, unit variance
    
    \item \textbf{Learning rate is key}: Spend time tuning it (try powers of 10)
    
    \item \textbf{Batch normalize}: Often helps training significantly
    
    \item \textbf{Use validation set}: Never tune hyperparameters on test set
    
    \item \textbf{Be patient}: Training large models takes time
    
    \item \textbf{Reproducibility}: Set random seeds for debugging
\end{enumerate}
\end{summary}

% ------------------------------------------------------------------------------
\section{From Scratch to Frameworks}
% ------------------------------------------------------------------------------

The autograd engine we built demonstrates the core concepts. Real frameworks add:

\begin{itemize}
    \item \textbf{GPU acceleration}: Massive parallelism for matrix operations
    \item \textbf{Optimized kernels}: Highly tuned CUDA/cuDNN implementations
    \item \textbf{Advanced optimizers}: Adam, RMSprop, AdaGrad, etc.
    \item \textbf{Built-in layers}: Conv2D, BatchNorm, Transformer, etc.
    \item \textbf{Data pipelines}: Efficient loading and augmentation
    \item \textbf{Distributed training}: Scale across multiple machines
\end{itemize}

\begin{lstlisting}[language=Python, caption={Same Network in PyTorch}]
import torch
import torch.nn as nn

# Define network
model = nn.Sequential(
    nn.Linear(3, 4),
    nn.Tanh(),
    nn.Linear(4, 4),
    nn.Tanh(),
    nn.Linear(4, 1)
)

# Data
X = torch.tensor([[2.0, 3.0, -1.0],
                  [3.0, -1.0, 0.5],
                  [0.5, 1.0, 1.0],
                  [1.0, 1.0, -1.0]])
y = torch.tensor([[1.0], [-1.0], [-1.0], [1.0]])

# Training
optimizer = torch.optim.SGD(model.parameters(), lr=0.1)
criterion = nn.MSELoss()

for epoch in range(100):
    pred = model(X)
    loss = criterion(pred, y)
    
    optimizer.zero_grad()
    loss.backward()
    optimizer.step()
    
    if epoch % 10 == 0:
        print(f"Epoch {epoch}, Loss: {loss.item():.4f}")
\end{lstlisting}

The structure is identical to our from-scratch implementation, but now running efficiently on GPU with production-ready code!

% ------------------------------------------------------------------------------
\section{Key Takeaways}
% ------------------------------------------------------------------------------

\begin{summary}
\textbf{Core Concepts:}
\begin{itemize}
    \item \textbf{Automatic differentiation} (autograd) efficiently computes gradients by recording operations and applying the chain rule
    \item \textbf{Reverse-mode autograd} is exactly backpropagation
    \item Modern frameworks (PyTorch, TensorFlow) are sophisticated autograd engines with GPU support
\end{itemize}

\textbf{The Training Loop:}
\begin{enumerate}
    \item Forward pass (compute predictions)
    \item Compute loss
    \item Zero gradients
    \item Backward pass (compute gradients)
    \item Update parameters
    \item Repeat
\end{enumerate}

\textbf{Key Challenges:}
\begin{itemize}
    \item Vanishing/exploding gradients $\rightarrow$ proper initialization, ReLU, normalization
    \item Overfitting $\rightarrow$ regularization, more data, simpler models
    \item Hyperparameter tuning $\rightarrow$ learning rate most critical
\end{itemize}
\end{summary}
